\chapter{Objectives and Investigation Question}
\label{chap:final-objectives}

\section{General Objective}

The general objective of this dissertation is to develop, evaluate and rigorously compare a\acrshort{cnn} model for metal artifact reduction in\acrshort{ct}, prioritizing anatomical preservation, robustness to domain variation and interpretability of the correction process. Additionally, the study aims to investigate the generalization of the best performing model to real and external data, with a focus on clinical applicability. Healthcare professionals will be involved in the evaluation of the model's performance on real data, providing a comprehensive assessment of its effectiveness and potential for clinical use.

\section{Specific Objectives}
The specific objectives of this study include:

\begin{itemize}
    \item Develop a reproducible pipeline for training and evaluating metal artifact reduction models in\acrshort{ct} in the image domain, with rigorous separation between training, validation, and test data.

    \item Evaluate the performance of a\acrfull{unet} for image reconstruction in metal artifact reduction and compare it with a version of itself with an additional artifact-mask supervision output.
    \item Evaluate the performance of the two versions using quantitative metrics like: \acrfull{psnr},\acrfull{ssim},\acrfull{mae}, \acrfull{bce} and \acrfull{dice}. Besides, performing error maps, statistical and qualitative analysis to assess the quality of the reconstructed images.
    \item Investigate the generalization capability of the best performing version in the internal evaluation stage on real and external data from the\acrfull{tcia}.
\end{itemize}

\section{Hypothesis and Outcomes}

The primary hypothesis of this study is that, in the internal test partition, the correction produced by each architecture makes a real, meaningful difference to the structural quality of the image compared to the artifact input. This is measured with\acrshort{ssim}, calculated separately in the\acrfull{roi} where the metal artifact was inserted and in the clean\acrshort{roi}, without an artifact. The clean\acrshort{roi} should stay about the same or change only a little. Each patient's slices were averaged into a single value per patient ($n=41$), and the comparison asked whether there was typically no difference at all between the input and the model.

The secondary hypotheses are threefold. First, that\acrshort{psnr} and\acrshort{mae}, in HU, also improve compared to the input, using the same per-patient averaging. Second, that the same three metrics (\acrshort{ssim},\acrshort{psnr} and\acrshort{mae}) also improve when calculated over the whole image. Third, a direct comparison between the Single Model and the Dual Model on these same three metrics, addressing the architecture-comparison objective.

As a way to judge anatomical preservation, the change introduced in the Clean Zone was expected to be clearly smaller than the correction observed in the Artifact Zone. The overall shift in average HU values, checked with a Bland--Altman plot at the patient level, was expected to be small compared to the width of the\acrshort{hu} window used. No fixed numeric limit was set in advance for how small the Clean-Zone change had to be, so preservation is described rather than formally tested.

Finally, since the external\acrshort{tcia} set has no clean reference image to compare against, the generalization hypothesis is necessarily more limited. The selected architecture is expected to produce changes of a plausible size in the external cases, without simply reproducing the input unchanged. However, only the doctors' blinded evaluation (\autoref{sec:avaliacao-medicos}), not a reference-based metric, can confirm whether artifact removal is accurate and clinically safe.

